%%%%%%%%%%%%%%%%%%%%%%%%%%%%%%%%%%%%%%%%%%%%%%%%%%%%%%%%%%%%%%%%%%%%%%%%%%%%%%%%%%%%%%
%%%%%%%%%%%%%%%%%%%%  INGRESAR LAS PREGUNTAS DE LA  EVALUACIÓN %%%%%%%%%%%%%%%%%%%%%%%%
%%%%%%%%%%%%%%%%%%%%%%%%%%%%%%%%%%%%%%%%%%%%%%%%%%%%%%%%%%%%%%%%%%%%%%%%%%%%%%%%%%%%%%%

\paragraph{I} Read the following questions and identify the option that best answers each, then write ts letter on the line (5 points each)

\begin{preguntas}

\pregunta[5] \rule{1cm}{0.1mm} Select the option with a quadratic function:

\begin{enumerate*}[label=(\alph*)]
    \item $(a/b)x + c\qquad$
    \item $(a_1/b_1)x + (a_2/b_2)x^2\qquad$
    \item $(a/b)x^3 + c\qquad$
    \item $(a_n/b_n)x^{1/2}\qquad$
\end{enumerate*}

\pregunta[5] \rule{1cm}{0.1mm} Select the correct domain, such that $f(x)\in\mathbb{R}$, where $f(x):=x^{1/n}$

\begin{enumerate*}[label=(\alph*)]
    \item $x\in\left(-\infty,\infty\right)\qquad$
    \item $x\in\left[0,\infty\right)\qquad$
    \item $x\in\left(-\infty,0\right]\qquad$
    \item $x\in\left[-1,1\right]\qquad$
\end{enumerate*}

\pregunta[5] \rule{1cm}{0.1mm} Select the option with the range of the function $f(x):=\mathrm{abs}(x)$ for $x\in\left(-\infty,\infty\right)$

\begin{enumerate*}[label=(\alph*)]
    \item $f(x)\in\left(-\infty,\infty\right)\qquad$
    \item $f(x)\in\left[0,\infty\right)\qquad$
    \item $f(x)\in\left(-\infty,0\right]\qquad$
    \item $f(x)\in\left[-1,1\right]\qquad$
\end{enumerate*}

\pregunta[5] \rule{1cm}{0.1mm} Select the correct type of function of $f(x):=x^{1/n}$, where $n\in\mathbb{Z}$ ($x$ belongs to the set of integer numbers)

\begin{enumerate*}[label=(\alph*)]
    \item Quadratic$\qquad$
    \item Cubic$\qquad$
    \item Exponential$\qquad$
    \item Rational$\qquad$
\end{enumerate*}

\pregunta[5] \rule{1cm}{0.1mm} Select the alternative mathematical notation to express $g[f(x)]$

\begin{enumerate*}[label=(\alph*)]
    \item $f\circ g\qquad$
    \item $g\circ f\qquad$
    \item $f\to g\qquad$
    \item $f(g)\qquad$
\end{enumerate*}

\pregunta[5] \rule{1cm}{0.1mm} Select the option of the function that map all the domain into the positive real numbers.

\begin{enumerate*}[label=(\alph*)]
    \item $f(x):=x^3\qquad$
    \item $f(x):=\mathrm{abs}(x)x\qquad$
    \item $f(x):=\mathrm{sgn}(x) x\qquad$
    \item $f(x):=-x^2\qquad$
\end{enumerate*}


\end{preguntas}

\newpage

\paragraph{II} Solve the following exercises in an orderly and clear manner. 
\textbf{Answers without procedure will not be considered.}
Underline or frame your final answer.

\begin{preguntas}

\pregunta[12] Tacking into account the following functions, 
\begin{gather*}
    f(x):=\frac{a}{b}x+c,\quad g(x):=x+h,
\end{gather*}
express the algebraic representation of 
\begin{gather*}
    H(x):= \frac{(f\circ g)(x) - f(x)}{h}.
\end{gather*}

\begin{center}
    \framebox(0.9\textwidth,200){}   
\end{center}


    

\pregunta[12] Tacking into account the following functions,
\begin{gather*}
    f(x):=\frac{a}{b}x,\quad g(x):=\frac{c}{d}x^2,
\end{gather*}
represent the algebraic expression of,
\begin{gather*}
    \frac{f(x)}{g(x)},\quad \frac{g(x)}{f(x)}.
\end{gather*}
Finally, write the main difference between them.

\begin{center}
    \framebox(0.9\textwidth,200){}   
\end{center}

\newpage

\pregunta[12] Use the verbal representation of the following functions,

\begin{enumerate*}[label=(\alph*)]
    \item $f(x):= a x + c\qquad$
    \item $f(x):= \left(a^2 + b^2\right)^{1/2}\qquad$
    \item $f(x):= \left(x-1\right)^{1/2},x\in(1,10]\qquad$
\end{enumerate*}

\begin{center}
    \framebox(0.9\textwidth,200){}   
\end{center}


\pregunta[12] Write a general algebraic expression for each graph,

\begin{figure}[ht!]
    \centering
    \begin{subfigure}[c]{0.45\textwidth}
        \centering   
        \includegraphics[width=0.75\textwidth]{NO TOCAR/imagenes/inverse-cuadratic.png}
    \end{subfigure}
    \begin{subfigure}[c]{0.45\textwidth}
         \centering
         \includegraphics[width=0.75\textwidth]{NO TOCAR/imagenes/cubic.png}
    \end{subfigure}

    \vspace{1.5cm}

    \begin{subfigure}[c]{0.45\textwidth}
         \centering
         \includegraphics[width=0.75\textwidth]{NO TOCAR/imagenes/linear.png}
    \end{subfigure}
    \begin{subfigure}[c]{0.45\textwidth}
        \centering
        \includegraphics[width=0.75\textwidth]{NO TOCAR/imagenes/inverse-cubic.png}
    \end{subfigure}
\end{figure}

\newpage

\pregunta[12] Graph a sketch of the following function,
\begin{gather*}
    f(x):=\left\{
    \begin{array}{cc}
        x^2 & x\in[0,\infty) \\
        -x^2 & x\in(-\infty,0] 
    \end{array}
    \right.
\end{gather*}

\begin{figure}[ht!]
    \centering
    \includegraphics[width=0.7\textwidth]{NO TOCAR/imagenes/cartesian-plane.png}
\end{figure}


\pregunta[10] Express the domain and range of the following functions,

\begin{enumerate*}[label=(\alph*)]
    \item $f(x):= (x+1)^{1/2}\qquad$
    \item $g(x):= 1/x^2\qquad$
\end{enumerate*}

\begin{center}
    \framebox(0.9\textwidth,200){}   
\end{center}


\end{preguntas}


\begin{comment}

\begin{preguntas}

%%%%%%%%%%%%%%%%  Pregunta %%%%%%%%%%%%%%%%
\pregunta[5] What is function?

%%%%%%%%%%%%%%%%  Pregunta %%%%%%%%%%%%%%%%
\pregunta[5] What is the difference between Domain and Range.


%%%%%%%%%%%%%%%%  Pregunta %%%%%%%%%%%%%%%%
\pregunta[5] Represent the following function in a verbal format,
\begin{equation*}
    f(x):=\frac{a_n}{b_n}x^n+c_n.
\end{equation*}

%%%%%%%%%%%%%%%%  Pregunta %%%%%%%%%%%%%%%%
\pregunta[5] What are the 4 ways to represent a function?


%%%%%%%%%%%%%%%%  Pregunta %%%%%%%%%%%%%%%%
\pregunta[5] Use mathematical notation to represent the following: ``$x$ has a value between $0$ and $16$''.


%%%%%%%%%%%%%%%%  Pregunta %%%%%%%%%%%%%%%%
\pregunta[5] What is the difference between exponential and rational functions?

\end{preguntas}

\begin{preguntas}
    
\pregunta[4] Calcule el siguiente límite:
\[
\lim_{x \to 2} \frac{x^2 - 4}{x - 2}
\]
Explique su procedimiento paso a paso.



%%%%%%%%%%%%%%%%  Pregunta 2 %%%%%%%%%%%%%%%%

\pregunta[4] Determine si la siguiente función es continua en $x = 3$. Justifique su respuesta.
\[
f(x) = 
\begin{cases} 
x^2 + 1 & \text{si } x < 3 \\
7 & \text{si } x = 3 \\
2x + 1 & \text{si } x > 3 
\end{cases}
\]



%%%%%%%%%%%%%%%%  Pregunta 3 %%%%%%%%%%%%%%%%

\pregunta[4] Calcule la derivada de las siguientes funciones:
             \begin{partes}
                 \parte[2] Pregunta a
                 \parte[2] Pregunta b
             \end{partes}



%%%%%%%%%%%%%%%%  Pregunta 4 %%%%%%%%%%%%%%%%

\pregunta[4] Utilice la regla del producto para calcular la derivada de la siguiente función:
\[
f(x) = (2x^2 + 3)(x^3 - 4x)
\]
Explique cada paso del procedimiento. 
    \fuente{Adaptado de Curo, 2014}


%%%%%%%%%%%%%%%%  Pregunta 5 %%%%%%%%%%%%%%%%

\pregunta[4] Halle la segunda derivada de la función:
\[
f(x) = \frac{x^4}{4} - 2x^2 + 5
\]
Muestre los pasos necesarios.

\end{preguntas} 

\end{comment}
